% TOP_PROBLEM: {"id": 1515, "title": "Navier-Stokes Existence and Smoothness", "category_id": 9, "category": {"id": 9, "name": "pde", "display_name": "Partial Differential Equations", "description": "PDEs and their applications in physics and geometry.", "slug": "pde", "order_index": 9, "created_at": "2026-07-31T15:26:25.671Z"}, "status": "open"}
% FIRST_POSTED: null
% ENTRY_KIND: statement_only
% Imported from: batch08.tex; UnsolvedMath ID 1515
\documentclass[11pt]{article}
\usepackage[margin=25mm]{geometry}
\usepackage{fontspec}
\setmainfont{Latin Modern Roman}
\usepackage{amsmath,amssymb,mathtools}
\usepackage{unicode-math}
\setmathfont{Latin Modern Math}
\usepackage{xurl,hyperref}
\hypersetup{colorlinks=true,urlcolor=blue}
\setcounter{secnumdepth}{0}
\setlength{\parindent}{0pt}
\setlength{\parskip}{5pt}
\setlength{\emergencystretch}{3em}
\newcommand{\sref}[2]{\hyperlink{source:#1:#2}{[S#2]}}
\newcommand{\cataloguescope}[1]{\paragraph{Recorded target.}#1}
\begin{document}
% UnsolvedMath ID 1515; source catalogue code PDE-001
\setcounter{section}{1514}
\section{Navier-Stokes Existence and Smoothness}\label{1515}
{\small\sffamily UnsolvedMath ID 1515 · Partial Differential Equations}
\subsection{Definitions and mathematical statement}

\paragraph{Shared notation.}$\mathbb N=\{1,2,\ldots\}$, and $\mathbb Z,\mathbb Q,\mathbb R,\mathbb C$ denote the integers, rationals, reals and complex numbers. Any other conventions are specified locally.


Fix a viscosity \(\nu>0\). For \(x\in\mathbb R^3\) and \(t\ge0\),
the unknown velocity \(u=(u_1,u_2,u_3)\) and pressure \(p\)
obey the incompressible Navier--Stokes equations
\[
 \partial_tu+(u\cdot\nabla)u=\nu\Delta u-\nabla p+f,
 \qquad \nabla\cdot u=0,
 \qquad u(x,0)=u^0(x).
\]
Here \(((u\cdot\nabla)u)_i=\sum_{j=1}^3u_j\partial_{x_j}u_i\)
and \(\Delta=\sum_{j=1}^3\partial_{x_j}^2\).
The initial velocity \(u^0\) is smooth and divergence free; \(f\)
is a prescribed smooth external force.

In the Euclidean formulation, admissible data and forces satisfy,
for every multi-index \(\alpha\in\mathbb Z_{\ge0}^3\) and all integers
\(m,K\ge0\), bounds with finite constants
\[
 |\partial_x^\alpha u^0(x)|\le C_{\alpha K}(1+|x|)^{-K},\qquad
 |\partial_x^\alpha\partial_t^mf(x,t)|
       \le C_{\alpha mK}(1+|x|+t)^{-K}.
\]
A physically admissible Euclidean solution has
\(u,p\in C^\infty(\mathbb R^3\times[0,\infty))\) and a bound
\[
                     \sup_{t\ge0}\int_{\mathbb R^3}|u(x,t)|^2\,dx<\infty.
\]

In the periodic formulation, \(u^0\) and \(f\) are periodic
under \(x\mapsto x+e_j\), where \(e_j\) are the three standard unit
vectors, so the spatial domain is \(\mathbb T^3=\mathbb R^3/\mathbb Z^3\).
The initial velocity is smooth and divergence free, and the force satisfies
\[
 |\partial_x^\alpha\partial_t^mf(x,t)|
       \le C_{\alpha mK}(1+t)^{-K}
                       \quad\text{for every }\alpha,m,K.
\]
An admissible periodic solution has \(u,p\) smooth for all \(t\ge0\)
and periodic in each spatial coordinate, including the pressure as required
by the official erratum.
\paragraph{Statement recorded in the supplied source.}
Prove at least one of the four official alternatives:
\begin{enumerate}
 \item[(A)] For every admissible Euclidean \(u^0\), with \(f\equiv0\),
 there exists a global physically admissible Euclidean solution.
 \item[(B)] For every admissible periodic \(u^0\), with \(f\equiv0\),
 there exists a global admissible periodic solution.
 \item[(C)] There exist admissible Euclidean \(u^0\) and \(f\)
 for which no global physically admissible Euclidean solution exists.
 \item[(D)] There exist admissible periodic \(u^0\) and \(f\)
 for which no global admissible periodic solution exists.
\end{enumerate}
In all four alternatives the viscosity is positive; existence alternatives
are unforced, whereas the breakdown alternatives permit the specified force. \sref{1515}{2}
\subsection{Short English statement}
Do smooth incompressible three-dimensional fluid data always yield global smooth solutions, or can admissible data and forcing prevent them? The official question permits an answer in Euclidean space or on the periodic three-torus.
\subsection{Sources}
\begin{description}
\item[\hypertarget{source:1515:1}{S1}] ULAM.AI and the UnsolvedMath Contributors, \emph{UnsolvedMath: A Curated Collection of Open Mathematics Problems}, record 1515 (PDE-001). CC BY 4.0. \url{https://huggingface.co/datasets/ulamai/UnsolvedMath}.
\item[\hypertarget{source:1515:2}{S2}] Charles L. Fefferman, \emph{Existence and Smoothness of the Navier--Stokes Equation}, official Clay Millennium problem description, pp. 1--2, equations (1)--(11), alternatives (A)--(D), and appended pressure-periodicity erratum. \url{https://www.claymath.org/wp-content/uploads/2022/06/navierstokes.pdf}.
\end{description}
\label{end:1515}
\end{document}
